% Edited conversion of source pp. 33--34.
\ReportHeading
  {Роль SH-хвиль у дослідженні хвильової динаміки шаруватих середовищ}
  {А.~Ю. Глухов}
  {Інститут механіки ім. С.~П. Тимошенка НАН України, Київ, Україна}
  {}

Розглянуто хвилі кручення, або SH-хвилі (\emph{Shear Horizontal}), у багатошарових середовищах із початковими напруженнями. Поряд із поздовжніми P-хвилями та вертикально поляризованими зсувними SV-хвилями вони належать до фундаментальних типів пружних хвиль.

SH-хвилі описуються скалярним хвильовим рівнянням. Рівняння руху допускає розділення змінних, що дає змогу одержувати аналітичні розв’язки у замкненій формі та прозоро досліджувати вплив шаруватості, граничних умов і початкових напружень. На межах поділу SH-хвилі не створюють нормальних напружень, тому умови спряження шарів істотно спрощуються. Оскільки ці хвилі є чисто зсувними, вони безпосередньо реагують на зміну модуля зсуву матеріалу.

Завдяки цим властивостям SH-хвилі є зручним модельним класом для аналізу хвильової динаміки шаруватих середовищ. Водночас конкретні розв’язки для них не можна безпосередньо переносити на всі типи хвиль. Узагальнюються насамперед установлені фізичні механізми: зміна швидкості поширення під дією початкових напружень, інтерференція внаслідок багаторазових відбиттів на межах шарів і локалізація напружень в окремих шарах.

Для P- і SV-хвиль виникає зв’язана система рівнянь, що описує взаємодію поздовжніх і поперечних деформацій. Відбиття та проходження через межі супроводжуються взаємним перетворенням хвиль, структура поля ускладнюється, а повне аналітичне розділення змінних, як правило, стає неможливим. Проте основні фізичні закономірності, виявлені для SH-хвиль, зберігаються і в загальніших хвильових режимах, хоча проявляються у складнішій формі.

Отже, дослідження SH-хвиль доцільно розглядати як базовий етап: воно поєднує можливість одержання аналітичних розв’язків із виявленням фізично значущих ефектів і створює основу для подальшого аналізу складніших задач хвильової динаміки багатошарових середовищ із початковими напруженнями.
